\documentclass[tikz,border=8pt]{standalone}
\usepackage{ctex}
\usepackage{amsmath}
\usepackage{amssymb}
\usepackage{pgfplots}
\pgfplotsset{compat=1.18}
\usetikzlibrary{calc,positioning,arrows.meta,decorations.pathreplacing}

\begin{document}
\begin{tikzpicture}

%% ===== 共轭梯度方向 vs 梯度下降方向 =====
%% 左图：普通梯度下降（"之字"路径）
%% 右图：共轭梯度方向（直接对角穿越椭圆）
%% 椭圆等高线：f(x,y) = x^2/9 + y^2  (条件数 kappa = 9)
%% 最优解 x* = (0, 0)

%% ========================
%% 左图：梯度下降（之字形）
%% ========================
\begin{scope}[xshift=0cm]

  %% 坐标轴
  \draw[->, gray!40, thin] (-3.8, 0) -- (3.8, 0)
    node[right, font=\scriptsize, gray] {$x_1$};
  \draw[->, gray!40, thin] (0, -2.4) -- (0, 2.4)
    node[above, font=\scriptsize, gray] {$x_2$};

  %% 等高线（椭圆，半轴比 3:1）
  \foreach \c in {0.3, 0.7, 1.2, 1.8, 2.4} {
    \draw[gray!35, thin] (0,0) ellipse ({3*\c} and {\c});
  }

  %% 极小值点
  \fill[green!50!black] (0,0) circle (2.5pt);
  \node[below right, font=\scriptsize, green!50!black] at (0.1,-0.1) {$\mathbf{x}^*$};

  %% 初始点
  \coordinate (GD0) at (-3.0, 1.8);
  \fill[black] (GD0) circle (2.5pt);
  \node[above left, font=\scriptsize] at (GD0) {$\mathbf{x}_0$};

  %% 梯度下降路径（之字形）
  %% 梯度 = (2x/9, 2y)，在 (-3, 1.8) 处：(-0.667, 3.6)，方向 (-1, 5.4) 归一化
  %% 沿负梯度精确线搜索后到达正交于下一个梯度的点
  %% 手工计算简化路径，模拟典型之字
  \coordinate (GD1) at (2.7, 0.3);   %% 第一步后（近似精确线搜索）
  \coordinate (GD2) at (-2.4, -0.25);  %% 第二步
  \coordinate (GD3) at (2.1, 0.06);   %% 第三步
  \coordinate (GD4) at (-1.8, -0.04); %% 第四步（仍未到达）

  \draw[->, thick, red!70!black, >=Stealth]
    (GD0) -- (GD1) node[midway, above, font=\scriptsize, red!70!black] {$-\nabla f_0$};
  \draw[->, thick, red!70!black, >=Stealth]
    (GD1) -- (GD2) node[midway, below, font=\scriptsize, red!70!black] {$-\nabla f_1$};
  \draw[->, thick, red!70!black, >=Stealth]
    (GD2) -- (GD3) node[midway, above, font=\scriptsize, red!70!black] {$-\nabla f_2$};
  \draw[->, thick, red!70!black, dashed, >=Stealth]
    (GD3) -- (GD4) node[midway, below, font=\scriptsize] {$\cdots$};

  \fill[red!70!black] (GD1) circle (2pt);
  \fill[red!70!black] (GD2) circle (2pt);
  \fill[red!70!black] (GD3) circle (2pt);

  %% 标注
  \node[font=\small\bfseries, red!70!black] at (0, -2.7)
    {梯度下降：沿 $-\nabla f$（"之"字振荡）};

\end{scope}

%% ========================
%% 右图：共轭梯度（直线穿越）
%% ========================
\begin{scope}[xshift=8.5cm]

  %% 坐标轴
  \draw[->, gray!40, thin] (-3.8, 0) -- (3.8, 0)
    node[right, font=\scriptsize, gray] {$x_1$};
  \draw[->, gray!40, thin] (0, -2.4) -- (0, 2.4)
    node[above, font=\scriptsize, gray] {$x_2$};

  %% 等高线（同上）
  \foreach \c in {0.3, 0.7, 1.2, 1.8, 2.4} {
    \draw[gray!35, thin] (0,0) ellipse ({3*\c} and {\c});
  }

  %% 极小值点
  \fill[green!50!black] (0,0) circle (2.5pt);
  \node[below right, font=\scriptsize, green!50!black] at (0.1,-0.1) {$\mathbf{x}^*$};

  %% 初始点（同位置）
  \coordinate (CG0) at (-3.0, 1.8);
  \fill[black] (CG0) circle (2.5pt);
  \node[above left, font=\scriptsize] at (CG0) {$\mathbf{x}_0$};

  %% CG 第一步：沿 d_0 = -∇f_0（与 GD 相同）
  %% 精确线搜索，到达 x_1，此时 ∇f(x_1) ⊥ d_0（A-共轭等高线切点）
  \coordinate (CG1) at (2.7, 0.3);
  \fill[blue!70!black] (CG1) circle (2.5pt);
  \node[below right, font=\scriptsize, blue!70!black] at (CG1) {$\mathbf{x}_1$};

  %% CG 第二步：d_1 = -∇f_1 + beta_1 * d_0
  %% beta 使 d_1 与 d_0 共轭（A-正交），方向斜着指向 x*
  %% 第二步直接到达 x*（n=2 维精确终止）
  \draw[->, thick, blue!70!black, >=Stealth]
    (CG0) -- (CG1)
    node[midway, above, font=\scriptsize, blue!70!black] {$\mathbf{d}_0 = -\nabla f_0$};
  \draw[->, thick, blue!70!black, >=Stealth]
    (CG1) -- (0, 0)
    node[midway, below right, font=\scriptsize, blue!70!black]
    {$\mathbf{d}_1 = -\nabla f_1 + \beta_1 \mathbf{d}_0$};

  %% 标注 A-共轭
  \draw[decorate, decoration={brace, amplitude=5pt, raise=2pt},
        gray!60]
    (CG0) -- (0, 0)
    node[midway, above, xshift=2pt, yshift=8pt, font=\scriptsize, gray!70!black]
    {2 步精确到达 $\mathbf{x}^*$（$n=2$）};

  %% beta 标注箭头
  \node[draw=blue!20, fill=blue!5, rounded corners=2pt,
        font=\scriptsize, inner sep=3pt, align=left]
    at (3.5, 2.0) {
      $\mathbf{d}_1 \perp_{\mathbf{A}} \mathbf{d}_0$\\
      （$\mathbf{A}$-共轭方向）\\
      $\beta_1 = \dfrac{\|\mathbf{r}_1\|^2}{\|\mathbf{r}_0\|^2}$
    };

  %% 标注
  \node[font=\small\bfseries, blue!70!black] at (0, -2.7)
    {共轭梯度：A-共轭方向，$n$ 步精确解};

\end{scope}

%% ===== 整体标题 =====
\node[font=\normalsize\bfseries] at (4.25, 3.3) {%
  共轭梯度方向（右）vs 梯度下降（左）：椭圆等高线上的路径对比};

\end{tikzpicture}
\end{document}
